% TÜBİTAK 2209-A araştırma önerisi. Derleme: tectonic oneri.tex  (veya xelatex oneri.tex)
\documentclass[11pt,a4paper]{article}
\usepackage[margin=2.5cm]{geometry}
\usepackage[shorthands=off,turkish]{babel}
\usepackage{amsmath}
\usepackage{array,booktabs,tabularx,longtable}
\usepackage[hidelinks]{hyperref}
\usepackage{xurl}
\usepackage{needspace}

\newcommand{\doi}[1]{DOI: \href{https://doi.org/#1}{\nolinkurl{#1}}}
\newcolumntype{L}[1]{>{\raggedright\arraybackslash}p{#1}}
\newcolumntype{Y}{>{\raggedright\arraybackslash}X}
\setlength{\parindent}{0pt}
\setlength{\extrarowheight}{3pt}
\setlength{\parskip}{0.5\baselineskip}

\begin{document}

\begin{center}
  {\Large\bfseries TÜBİTAK 2209-A Araştırma Önerisi --- 2026 Yılı Çağrısı\par}
\end{center}

\textbf{Proje Başlığı:} Kalın Kesitli PSMA PET/BT Görüntülerinde BT Kılavuzlu Derin Öğrenme ile
Kesitler Arası Çözünürlük Artırımı ve Lezyon Düzeyinde SUV Sadakatinin Korunması

\textbf{Anahtar Kelimeler:} PSMA PET/BT, kesit interpolasyonu, süper çözünürlük, derin öğrenme,
çapraz dikkat, kısmi hacim etkisi, SUV

\section*{Özet}

Prostata özgü membran antijeni (PSMA) PET/BT, prostat kanserinin evrelemesinde ve PSMA hedefli
radyoligand tedavisine hasta seçiminde klinik standart hâline gelmiştir (Hofman vd., 2020;
Fendler vd., 2023; Sartor vd., 2021). PET görüntülerinin eksenel (z) yöndeki kesit aralığı, aynı
seansta çekilen BT'ye göre belirgin biçimde kabadır; kalın kesitlerde küçük lezyonların aktivitesi
komşu dokuya yayılır ve kısmi hacim etkisi nedeniyle maksimum standart tutulum değeri (SUVmax)
olduğundan düşük ölçülür (Soret vd., 2007). Bu projede, seyrek (kalın) PET kesitlerinden ince
kesitli PET hacmini, aynı seansta elde edilen tam çözünürlüklü BT'nin anatomik ön bilgisi
kılavuzluğunda yeniden oluşturan bir derin öğrenme yöntemi geliştirilecektir. Problem, kalın
kesitin ince kesitlerin ortalaması olduğu fiziksel olarak tutarlı bir bozunum modeliyle
tanımlanacak ve $k = 2$, 3 ve 4 seyrekleştirme faktörlerinin her biri için ayrı ayrı
çözülecektir. Yöntem, kübik spline taban çizgisinden başlayıp BT'siz ve BT'li 2.5 boyutlu
rezidüel U-Net'e, önceden eğitilmiş bir BT temel modelinin özniteliklerine ve projenin özgün
katkısı olan, BT kodlayıcısından PET kod çözücüsüne çapraz dikkat (cross-attention) ile bilgi
aktaran mimariye uzanan bir model merdiveni olarak kurgulanmıştır. Tüm modellerin çıktısına,
sentezlenen ince kesitlerin ortalamasını ölçülen kalın kesite eşitleyen ölçüm tutarlılığı
izdüşümü uygulanacaktır. Literatürdeki SAINT, I$^3$Net, RIFE ve FILM yöntemleri aynı veri
bölünmesi, aynı bozunum modeli ve aynı metrik koduyla yeniden çalıştırılarak adil bir
karşılaştırma yapılacaktır. Değerlendirme, vücut maskesi içindeki PSNR, SSIM ve NRMSE'nin
yanında lezyon düzeyinde SUVmax sapmasını, 1~mL'den küçük lezyonları ve mesaneye komşu
lezyonları ayrı ayrı kapsayacak; karşılaştırmalar hasta düzeyinde eşleştirilmiş Wilcoxon
işaretli sıralar testi ve Bonferroni düzeltmesiyle yapılacaktır. Genellenebilirlik farklı bir
merkezin PSMA verisinde ve aynı hastaların FDG taramalarında sınanacaktır. Ön çalışmada,
$k = 2$'de öğrenmeli modelin kübik interpolasyonun \%13,2 olan ortalama lezyon SUVmax düşük
ölçümünü yaklaşık \%2'ye indirdiği, BT'nin girdi kanallarına basitçe eklenmesinin ise fayda
sağlamadığı görülmüştür; bu bulgu, BT bilgisinin ağa hangi yolla aktarılması gerektiği sorusunu
projenin merkezine yerleştirmektedir. Proje, üç lisans öğrencisinden oluşan ekip tarafından
akademik danışman rehberliğinde, 12 ay süren beş iş paketi hâlinde yürütülecektir. Her deney
sürümlenmiş yapılandırma dosyalarından yeniden üretilebilecek ve kesintiye uğradığında kaldığı
yerden sürdürülebilecek bir yazılım hattıyla gerçekleştirilecektir. Projenin çıktıları en az bir
uluslararası bildiri, bir dergi makalesi taslağı ve açık kaynak koddur. Uzun vadede yöntemin, PET
çekim süresini veya hasta dozunu artırmadan küçük lezyonların daha doğru ölçülmesine katkı
sağlaması hedeflenmektedir.

\section{Araştırma Önerisinin Bilimsel Niteliği}

\subsection{Konunun Önemi ve Araştırma Önerisinin Bilimsel Niteliği}

\paragraph{Klinik önem.} PSMA PET/BT, yüksek riskli prostat kanserinde konvansiyonel
görüntülemeye (BT ve kemik sintigrafisi) göre daha yüksek tanısal doğruluk sağladığı için
evrelemede klinik standart hâline gelmiştir (Hofman vd., 2020); görüntülerin elde edilmesi ve
yorumlanması uluslararası kılavuzlarla tanımlanmıştır (Fendler vd., 2023). Metastatik
kastrasyona dirençli prostat kanserinde $^{177}$Lu-PSMA-617 tedavisinin sağkalım yararı
gösterildikten sonra (Sartor vd., 2021) PSMA PET yalnızca bir tanı aracı değil, tedaviye hasta
seçimi ve yanıt değerlendirmesi için nicel bir ölçüm aracı da olmuştur. Bu kullanımda karar,
büyük ölçüde lezyonların standart tutulum değerlerine (SUV) dayanır; dolayısıyla SUV'deki
sistematik bir düşük ölçüm, doğrudan klinik bir hataya dönüşebilir.

\paragraph{Problem.} PET/BT taramasında BT ince kesitli ve yüksek uzaysal çözünürlüklü elde
edilirken PET'in eksenel (z) yöndeki örnekleme aralığı daha kabadır. Bu çalışmada kullanılacak
açık PSMA veri setinde dahi PET'in kesit aralığı hastadan hastaya 1,5~mm ile 5~mm arasında
değişmektedir. Kalın bir PET kesiti, tarayıcının o kalınlık boyunca topladığı sayımların
ortalamasıdır; birkaç milimetrelik bir lezyon bu ortalama içinde komşu dokuyla karışır ve kısmi
hacim etkisi nedeniyle SUVmax değeri gerçek değerinden düşük ölçülür (Soret vd., 2007). Etki en
çok, klinik açıdan en önemli olan küçük lezyonlarda ortaya çıkar: klinik bir çalışmada 4~mm
yerine 2~mm vokselle geriçatım yapılması SUVmax değerini ortalama \%32, 0,75~mL'den küçük
lezyonlarda ise \%44 artırmıştır (Koopman vd., 2015). Daha ince kesitler elde etmek ise ya daha
uzun çekim ya da daha yüksek doz anlamına gelir. PET görüntü iyileştirme literatürü gürültü
azaltma ve düzlem içi çözünürlük artırımı üzerinde yoğunlaşmıştır (Liu vd., 2021); derin
öğrenmeli PET süper çözünürlük çalışmaları SUV doğruluğunu değerlendirmeye yeni yeni başlamış
(Endo vd., 2025), BT kılavuzunda eksenel yöndeki çözünürlük kaybı ise görece az çalışılmıştır.

\paragraph{Literatürdeki boşluk.} Kesitler arası interpolasyon ve kesit sentezi problemi BT ve
MR için çalışılmıştır. SAINT, sagital ve koronal düzlemlerde öğrenilen süper çözünürlüğü
birleştirerek BT kesitleri sentezler (Peng vd., 2020); I$^3$Net, kesit içi ve kesitler arası
bilgiyi frekans alanında da öğrenerek bu alanda güncel en iyi sonuçları raporlamıştır (Song vd.,
2024); SMORE ise MR için harici yüksek çözünürlüklü veri gerektirmeyen öz-denetimli bir yaklaşım
önerir (Zhao vd., 2021). Video kare interpolasyonu yöntemleri RIFE (Huang vd., 2022) ve FILM
(Reda vd., 2022) de tıbbi kesit interpolasyonunda karşılaştırma yöntemi olarak sıkça
kullanılmaktadır. PET tarafında ise anatomik kılavuz fikri çoğunlukla beyin PET'inde MR
kılavuzlu düzlem içi süper çözünürlük olarak ele alınmıştır (Song vd., 2020). BT'den PET sentezi
üzerine yapılan çalışmalar (Ben-Cohen vd., 2019), BT anatomisinin tutulumun nerede olabileceğine
dair bilgi taşıdığını, ancak tutulumun şiddetini tek başına belirleyemediğini göstermektedir; bu
nedenle seyrek ölçülmüş PET ile BT'yi birlikte kullanan bir kurgu, yalnızca BT'den PET üretmeye
göre klinik olarak çok daha savunulabilirdir. Bilgimiz dâhilinde, aynı seansta elde edilen BT'yi
kılavuz olarak kullanan PSMA PET kesitler arası interpolasyonu sistematik olarak
çalışılmamıştır. Ayrıca mevcut kesit interpolasyonu çalışmalarının çoğu düşük çözünürlüklü
girdiyi kesit atlayarak (decimation) üretmektedir; oysa kalın kesitli PET kesit atlamaya değil,
kesit ortalamasına karşılık gelir ve bu iki problem fiziksel olarak farklıdır. Son olarak bu
çalışmaların değerlendirmesi çoğunlukla tüm görüntü üzerinden hesaplanan PSNR ve SSIM ile
yapılmakta, lezyon düzeyindeki SUV hataları bu ortalamaların içinde görünmez kalmaktadır.

\paragraph{Araştırma sorusu ve hipotezler.} Projenin araştırma sorusu şudur: Aynı seansta
çekilen tam çözünürlüklü BT, kalın kesitli PSMA PET'ten ince kesitlerin geri kazanılmasında,
özellikle küçük lezyonların SUVmax değerinin korunmasında, PET'in kendi içeriğinin ötesinde ek
bilgi sağlar mı; sağlıyorsa bu bilgi ağa hangi yolla aktarılmalıdır? Bu soru üç hipotezle
sınanacaktır. H1: Öğrenmeli modeller, kübik spline interpolasyonuna göre lezyon SUVmax sapmasını
ve görüntü hatasını hasta düzeyinde anlamlı olarak azaltır. H2: BT bilgisinin ayrı bir kodlayıcı
ve çapraz dikkat ile ya da önceden eğitilmiş bir BT temel modelinin öznitelikleriyle aktarılması,
BT'siz modele göre özellikle yüksek seyrekleştirme faktöründe ($k = 4$) ve küçük lezyonlarda
iyileşme sağlar; BT'nin girdi kanallarına basitçe eklenmesi (erken füzyon) ise sağlamaz. H3:
Ölçüm tutarlılığı izdüşümü, sentezlenen hacmin ölçülen veriyle tutarlı olmasını garanti ederek
global hata metriklerini iyileştirir; lezyon SUVmax üzerindeki etkisi ayrıca ölçülecektir.

\paragraph{Ulusal öncelikler ile ilişkisi.} Proje, TÜBİTAK öncelikli alanlar listesindeki
``Yapay zekâ'' ve ``Tıbbi cihaz teknolojileri'' alanlarıyla doğrudan ilişkilidir (TÜBİTAK, 2025).
On İkinci Kalkınma Planı'nda ``yapay zekâ tabanlı sağlık teknolojilerinin geliştirilmesi''
(madde 441.5), ``tıbbi görüntüleme cihazları alanında kullanılan kritik teknolojiler'' (madde
449.2), sağlık teknolojilerinin stratejik bir Ar-Ge alanı olarak sürdürülmesi (madde 556) ve
dijital sağlık teknolojilerinin kullanımının yaygınlaştırılması (madde 715) hedeflenmekte;
anonimleştirilmiş sağlık verilerinin Ar-Ge amaçlı ikincil kullanımı için altyapı oluşturulması
öngörülmektedir (madde 716.3) (Strateji ve Bütçe Başkanlığı, 2023). 2030 Sanayi ve Teknoloji
Stratejisi'nde yapay zekâ ve tıbbi cihaz, öncelikli yüksek teknoloji alanları arasında
sayılmakta ve yapay zekâ destekli teşhis sistemlerinin radyoloji görüntülerini analiz ederek
kanser gibi ciddi hastalıkları daha hızlı ve doğru biçimde tespit edebileceği vurgulanmaktadır
(Sanayi ve Teknoloji Bakanlığı, 2025). Bu proje, tamamen anonimleştirilmiş açık veriyle
geliştirilen, tıbbi görüntüleme cihazının donanımını değiştirmeden çıktısının nicel doğruluğunu
yazılım düzeyinde artıran yapay zekâ tabanlı bir sağlık teknolojisi bileşeni ürettiği için bu
maddelerle örtüşmektedir. Projenin bu alanlara yararı, nükleer tıp görüntülemesinde yerli bir
yöntem birikimi oluşturması ve yapay zekâ ile tıbbi görüntüleme kesişiminde çalışabilecek
nitelikli insan kaynağı yetiştirmesidir.

\subsection{Amaç ve Hedefler}

Projenin amacı, kalın kesitli PSMA PET görüntülerinden ince kesitli PET hacmini, aynı seansta elde
edilen BT'nin anatomik bilgisinden yararlanarak, lezyon düzeyindeki nicel değerleri (SUV)
koruyacak biçimde yeniden oluşturan, fiziksel olarak tutarlı ve yeniden üretilebilir bir derin
öğrenme yöntemi geliştirmek ve bu yöntemi literatürdeki yöntemlerle adil koşullarda
karşılaştırmaktır. Bu amaç doğrultusunda aşağıdaki ölçülebilir hedefler belirlenmiştir.

\paragraph{Hedef 1 --- Tekrarlanabilir veri ve bozunum hattı (İP1).} 100 hastalık açık PSMA
PET/BT kohortunun SUV dönüşümünün, kalite kontrolünün ve BT'nin PET ızgarasına yeniden
örneklenmesinin tamamlanması; hasta düzeyinde \%70/\%10/\%20 eğitim/doğrulama/test bölünmesinin
ve $k = 2$, 3, 4 için test bozunumlarının bir kez üretilip dondurulması. Ölçüt: kohortun tamamı
için kalite kontrol raporu, dışlanan her hasta için yazılı gerekçe ve dondurulmuş bölünme
dosyası.

\paragraph{Hedef 2 --- Adil karşılaştırma tablosu (İP1--İP2).} Kübik spline ve dört literatür
yönteminin (SAINT, I$^3$Net, RIFE, FILM) aynı bölünme, aynı ortalama tabanlı bozunum ve aynı
metrik koduyla $k = 2$, 3 ve 4'te değerlendirilmesi. Ölçüt: 5 yöntem $\times$ 3 faktörlük
tablonun eksiksiz doldurulması.

\paragraph{Hedef 3 --- Model merdiveni ve BT katkısının ölçülmesi (İP3).} BT'siz ve BT'li (erken
füzyon) 2.5B rezidüel U-Net ile BT temel modeli öznitelikli U-Net'in her faktörde en az üç farklı
rastgele tohumla eğitilmesi. Ölçüt: öğrenmeli modelin kübik interpolasyona göre test hastalarının
çoğunluğunda daha iyi olması ve farkın Bonferroni düzeltmesi sonrası $p < 0{,}05$ düzeyinde
anlamlı olması (H1); BT'li ve BT'siz modeller arasındaki farkın hasta düzeyinde test edilmesi
(H2'nin ilk ayağı).

\paragraph{Hedef 4 --- Özgün çapraz dikkat modeli (İP4).} BT'yi ayrı bir kodlayıcıdan çapraz
dikkat ile PET kod çözücüsüne aktaran modelin geliştirilmesi ve ablasyonlarla incelenmesi.
Ölçüt: $k = 4$'te BT'siz U-Net'e karşı hasta düzeyinde eşleştirilmiş testte PSNR'de anlamlı
üstünlük (H2) ve 1~mL'den küçük lezyonlarda SUVmax düşük ölçümünün kübik interpolasyona göre en
az \%40 azaltılması (ön çalışmada kübik interpolasyonun $k = 4$'teki düşük ölçümü \%37,6'dır).

\paragraph{Hedef 5 --- Genellenebilirlik ve klinik görev (İP5).} Modellerin farklı bir merkezden
gelen PSMA verisinde (dış test) ve aynı hastaların FDG taramalarında (tracer değişimi)
değerlendirilmesi; ayrıca hazır bir lezyon segmentasyon modelinin gerçek ve sentezlenmiş PET
üzerindeki başarısının karşılaştırılması. Ölçüt: dış test kümesindeki performans düşüşünün ve
segmentasyon Dice farkının raporlanması.

\paragraph{Hedef 6 --- Paylaşım.} Kodun ve deney yapılandırmalarının açık kaynak olarak
yayımlanması; en az bir uluslararası bildiri ve bir dergi makalesi taslağının hazırlanması.

\section{Yöntem}

\subsection{Araştırma tasarımı}
Çalışma, açık erişimli ve anonimleştirilmiş görüntüler üzerinde yürütülecek retrospektif,
hesaplamalı ve deneysel bir çalışmadır. Bağımsız değişkenler model türü (model merdiveninin
basamağı), BT'nin kullanılıp kullanılmadığı ve kullanılıyorsa aktarım biçimi (erken füzyon, temel
model öznitelikleri, çapraz dikkat), seyrekleştirme faktörü $k \in \{2, 3, 4\}$, ölçüm
tutarlılığı izdüşümünün uygulanıp uygulanmadığı ve rastgele tohumdur. Bağımlı değişkenler vücut
maskesi içinde hesaplanan PSNR, SSIM ve NRMSE; lezyon düzeyinde SUVmax ve SUVmean sapması; alt
gruplarda (1~mL'den küçük lezyonlar, mesaneye 2~cm'den yakın lezyonlar) SUVmax sapması ve klinik
görev değerlendirmesinde segmentasyon Dice skorudur. Her model her faktörde hem ham çıktısıyla hem
ölçüm tutarlılığı sonrası çıktısıyla raporlanacağından, tutarlılık ablasyonu her deneyle birlikte
elde edilecektir.

\subsection{Veri (İP1)}
Ana eğitim ve geliştirme kümesi, metastatik kastrasyona dirençli prostat kanseri hastalarına ait,
lezyon kontürleri (toplam tümör yükü, TTB) bulunan ve her hastada hem PSMA hem FDG PET/BT içeren
DEEP-PSMA veri setidir (100 hasta; Jackson vd., 2025); lezyon kontürleri Ferdinandus vd. (2020) ve
Buteau vd. (2022) tarafından tanımlanan yöntemle oluşturulmuştur. Dış test için farklı bir
merkezde çekilmiş, elle işaretlenmiş lezyonlar içeren PSMA-PET-CT-Lesions koleksiyonu (Jeblick
vd., 2026) kullanılacaktır; bu koleksiyon, ABD Ulusal Sağlık Enstitüleri (NIH) tarafından
desteklenen The Cancer Imaging Archive (TCIA) deposunda yayımlanmaktadır (Clark vd., 2013).
Tracer genellenebilirliği DEEP-PSMA'daki aynı hastaların FDG taramaları üzerinde, gerektiğinde
açık FDG PET/BT veri setiyle (Gatidis vd., 2022) sınanacaktır. Bütün veri setleri
anonimleştirilmiş ve araştırma amaçlı kullanım lisansıyla yayımlanmıştır. Kurumsal (hastane)
verisi ancak etik kurul onayı alındıktan sonra yalnızca dış doğrulama amacıyla kullanılacaktır.

\subsection{Ön işleme ve kalite kontrolü (İP1)}
PET voksel değerleri vücut ağırlığına göre SUV'ye çevrilecektir. Her hasta bölünmeye girmeden
önce bir kalite kontrol kapısından geçecek; PET/BT yön tutarsızlığı, sayısal olmayan veya negatif
SUV ve 40'tan az kesit dışlama nedenidir ve her dışlama gerekçesiyle kaydedilecektir. SUV
dönüşümü, FDG taramalarında karaciğer SUVmean değerinin fizyolojik aralıkta (1,3--3,0) olmasıyla
doğrulanacaktır (Boellaard vd., 2015). PSMA'da karaciğer tutulumu tümör yükü arttıkça düştüğü
için (tümör ``sink'' etkisi; Gafita vd., 2022) karaciğer değeri kaydedilecek ancak eleme ölçütü
olarak kullanılmayacak, PSMA taraması aynı hastanın FDG taraması üzerinden kontrol edilecek ve
böbrek ile mesanenin en parlak yapılar olması ayrıca denetlenecektir. BT, PET ızgarasına yeniden
örneklenecek (tersi değil, çünkü düşük çözünürlüklü PET'i büyütmek yapay bilgi üretir) ve aynı
seansta elde edildiği için deformable kayıt yapılmayacaktır. BT değerleri $[-1000, +1000]$~HU
aralığında kırpılıp $[-1, 1]$ aralığına, PET ise $\log(1 + \mathrm{SUV})/\log(1 + 50)$
dönüşümüyle ölçeklenecektir; hasta bazında min-maks normalizasyon, SUV'nin mutlak anlamını
bozacağı için kullanılmayacaktır. Bütün metrikler ters dönüşümle SUV uzayında hesaplanacaktır.
Vücut maskesi BT'de $-500$~HU eşiği ve en büyük bağlı bileşenle, mesane maskesi ise
TotalSegmentator ile (Wasserthal vd., 2023) elde edilecektir. Veri, komşu kesitlerin test
kümesine sızmasını önlemek için kesit düzeyinde değil hasta düzeyinde, toplam tümör hacminin
üçte birlik dilimlerine göre katmanlanarak ve sabit bir tohumla bölünecek; bölünme dosyası bir
kez yazıldıktan sonra değiştirilmeyecektir.

\subsection{Bozunum modeli (İP1)}
Kalın kesit, $k$ ardışık ince kesitin ortalaması olarak üretilecektir:
\[
  y_j = \frac{1}{k} \sum_{i=0}^{k-1} x_{jk+i}.
\]
Bu, tarayıcının kesit kalınlığı boyunca sayımları toplamasının bir modelidir; literatürde sık
kullanılan kesit atlama (decimation) ise fiziksel olarak farklı bir problemi tanımlar ve
kullanılmayacaktır. Hedef (yer gerçeği) tarayıcının ürettiği özgün ince kesitli hacimdir, yani
model gerçek bir çekimi geri kazanmayı öğrenecektir. Test kümesinin bozunumları her eğitimde
yeniden örneklenmeyecek, bir kez üretilip diske yazılacak ve dondurulacaktır.

\subsection{Model merdiveni (İP3--İP4)}
Literatürde kesit interpolasyonunda problem formülasyonunun mimariden daha belirleyici
olabildiği bilindiğinden, merdivenin her basamağında aynı formülasyon sabit tutulacak ve
formülasyon ayrıca ablate edilecektir. Ortalama tabanlı bozunumda hiçbir ince kesit ``eksik''
değildir; her ince kesit ölçülmüş bir kalın kesitin içindedir. Bu nedenle temel birim blok
olarak seçilmiştir: model bir kalın kesiti alır ve içindeki $k$ ince kesiti üretir. Basamak 0
kübik spline interpolasyonudur ve tüm metrik altyapısının uçtan uca testi olarak da
kullanılacaktır. Basamak 1, U-Net (Ronneberger vd., 2015) mimarisinin MONAI (Cardoso vd., 2022)
gerçeklemesine dayanan 2.5 boyutlu rezidüel bir ağdır: girdi olarak hedef kalın kesit ve her iki
yanındaki ikişer kalın PET kesiti, bloğu kapsayan BT kesitleri ve kübik tahmin verilir; ağ kübik
tahmine eklenecek düzeltmeyi öğrenir. Son katman sıfırla başlatıldığından eğitilmemiş ağ tam
olarak kübik taban çizgisine eşittir; böylece elde edilen her kazanç doğrudan öğrenmeye
atfedilebilir. Bu basamak BT'li (erken füzyon) ve BT'siz olarak iki kolda eğitilecektir.
Basamak 2'de BT kanalları yerine, 148.000 BT taraması üzerinde öz-denetimli olarak önceden
eğitilmiş CT-FM temel modelinin (Pai vd., 2025) dondurulmuş kodlayıcısından alınan çok ölçekli
öznitelikler, sıfırla başlatılmış $1 \times 1$ evrişimlerle U-Net kodlayıcısına eklenecektir;
temel modelin kendi çözünürlüğünde (yaklaşık 1~mm) beslenmesinin etkisi ayrıca sınanacaktır.
Basamak 3 projenin özgün katkısıdır: PET ve BT için iki ayrı kodlayıcı kullanılacak, kod
çözücünün her ölçeğinde PET öznitelikleri sorgu (query), BT öznitelikleri anahtar ve değer
(key/value) olacak şekilde çok başlıklı çapraz dikkat blokları (Vaswani vd., 2017)
yerleştirilecek, bellek gereksinimini sınırlamak için Swin UNETR'deki gibi pencereli dikkat
(Hatamizadeh vd., 2022) kullanılacaktır. Gerekçe, kanal birleştirmenin BT'nin keskin anatomik
sınır bilgisini derin katmanlara taşıyamaması, çapraz dikkatin ise PET'in ``nereye bakacağını''
anatomiden öğrenebilmesidir. Zaman kalması hâlinde Basamak 4 olarak kod çözücü, koordinat
koşullu örtük bir fonksiyona (LIIF/ArSSR yaklaşımı; Chen vd., 2021; Wu vd., 2023) çevrilerek
$k = 2$, 3, 4 ve kesirli faktörleri tek modelle çözen ölçekten bağımsız bir sürüm denenecektir.
Hayali SUV değerleri üretme riski klinik olarak savunulamayacağı için çekişmeli (GAN) kayıp
kullanılmayacaktır.

\subsection{Ölçüm tutarlılığı (tüm basamaklar)}
Her modelin çıktısı, sentezlenen ince kesitlerin $k$'lı ortalaması ölçülen kalın kesite tam
olarak eşit olacak biçimde izdüşürülecektir:
\[
  \hat{x} \leftarrow \hat{x} + U\bigl(y - A\hat{x}\bigr).
\]
Burada $A$, $k$'lı ortalama işlecini, $U$ ise her kalın kesit değerini $k$ ince kesite kopyalayan
büyütme işlecini gösterir. Bu yaklaşım, MR geriçatımında derin ağlara eklenen veri tutarlılığı
katmanlarıyla (Schlemper vd., 2018) aynı ilkeye dayanır ve modelin ölçülen veriyle çelişen bir
sonuç üretmesini yapısal olarak engeller.

\subsection{Literatür yöntemleriyle karşılaştırma (İP2)}
SAINT ve I$^3$Net yazarlarının açık kodları bu projenin veri yükleyicisine bağlanacak ve kendi
yayınlarındaki kesit atlama protokolü yerine bu projenin ortalama tabanlı bozunumu, aynı bölünme
ve aynı metrik koduyla yeniden eğitilecektir; aksi hâlde karşılaştırma adil olmaz. RIFE ve FILM
önce önceden eğitilmiş ağırlıklarla, eğitimsiz olarak ardışık PET kesit çiftlerine uygulanacak,
ardından aynı koşullarda ince ayar yapılacaktır. Tüm yöntemler tek bir değerlendirme arayüzünden
geçirilecek; girdi düşük çözünürlüklü hacim ve BT, çıktı yüksek çözünürlüklü hacimdir.

\subsection{Eğitim (İP3--İP4)}
Kayıp fonksiyonu, görüntü onarımı için önerilen $0{,}84 \cdot (1 - \text{MS-SSIM}) + 0{,}16 \cdot
L_1$ birleşimidir (Zhao vd., 2017; Wang vd., 2003); saf SSIM ailesi kayıpların kararsızlığı
nedeniyle $L_1$ terimi korunacak, az sayıdaki lezyon vokselinin toplam kayıp içinde
kaybolmaması için $L_1$ lezyon voksellerinde 5 kat ağırlıklandırılacaktır. Ağlar PyTorch (Paszke
vd., 2019) ile, AdamW (Loshchilov ve Hutter, 2019) optimize edicisi ve kosinüs öğrenme hızı
azaltımıyla (Loshchilov ve Hutter, 2017) eğitilecektir. Veri artırma yalnızca sağ-sol çevirme ve
düzlem içi döndürmeyle sınırlı tutulacak, problem z yönünde tanımlı olduğu için z yönünde artırma
yapılmayacaktır. Model seçimi doğrulama kümesi üzerinden yapılacak; test kümesi her dondurulmuş
tasarım için yalnızca bir kez okunacak ve hipotezler sonuçlar görülmeden önce deney günlüğüne
yazılacaktır. Tohum varyansını ayırmak için her kol her faktörde en az üç tohumla
eğitilecektir.

\subsection{Değerlendirme ve istatistiksel yöntemler (İP3--İP5)}
Global metrikler (PSNR, SSIM, NRMSE; Wang vd., 2004), havanın metrikleri yapay olarak
şişirmemesi için yalnızca vücut maskesi içinde hesaplanacaktır. Lezyonlar TTB kontürlerinin
26-komşuluklu bağlı bileşenleri olarak tanımlanacak; her lezyon için SUVmax sapması
\[
  \frac{\mathrm{SUVmax}_{\text{tahmin}} - \mathrm{SUVmax}_{\text{gerçek}}}{\mathrm{SUVmax}_{\text{gerçek}}}
\]
hesaplanacak, ortalama ve \%95 güven aralığıyla ve Bland-Altman
grafikleriyle (Bland ve Altman, 1986) raporlanacaktır. 1~mL'den küçük lezyonlar ve mesaneye
2~cm'den yakın lezyonlar ayrı alt gruplar olarak raporlanacaktır; ikincisi, PSMA'da çok parlak
olan mesanenin komşu kesitlere ``sızma'' hatasının test edildiği yerdir. Sonuçlar ayrıca
hastaların özgün kesit aralığına göre (1,5--5~mm) kırılacaktır. Klinik görev değerlendirmesinde
nnU-Net (Isensee vd., 2021) tabanlı bir lezyon segmentasyon modeli gerçek ve sentezlenmiş PET
üzerinde çalıştırılıp Dice ve lezyon tespit farkları ölçülecektir. Bütün karşılaştırmalar hasta
düzeyinde eşleştirilmiş Wilcoxon işaretli sıralar testiyle (Wilcoxon, 1945) yapılacak, çoklu
karşılaştırmalar Bonferroni yöntemiyle (Dunn, 1961) düzeltilecektir; birden çok tohumla eğitilen
kollar için önce her hastanın metrikleri tohumlar üzerinden ortalanacaktır. En kötü sonuç veren
beş vaka ayrıca görsel olarak incelenecektir.

\subsection{Tekrarlanabilirlik}
Her deney, yol, hiperparametre, tohum ve bölünme dosyasını içeren sürümlenmiş bir yapılandırma
dosyasından üretilecek; her koşu kendi yapılandırma kopyasını, kod sürümünü ve çalışma ağacının
temiz olup olmadığını arşivleyecektir. Bulut ortamındaki oturum kesintilerine karşı uzun aşamalar
hasta başına çıktı yazıp tamamlananları atlayarak, eğitim ise her dönemde (epoch) kontrol noktası
kaydederek kaldığı yerden sürdürülecektir. Tablolar ve şekiller elle değil, metrik dosyalarından
betiklerle üretilecektir.

\subsection{Ön çalışma ve fizibilite}
Yöntemin uygulanabilirliğini sınamak için DEEP-PSMA kohortunun 100 PSMA taraması ön işlenmiş,
tamamı kalite kontrol kapısından geçmiş ve hasta düzeyinde 70/9/21 bölünme dondurulmuştur.
$k = 2$'de yürütülen pilot deneyde 21 test hastasında (1.508 lezyon), ölçüm tutarlılığı
sonrasında kübik spline lezyon SUVmax değerini ortalama \%13,2, 1~mL'den küçük lezyonlarda
\%15,7 düşük ölçerken BT'siz 2.5B U-Net bu sapmayı sırasıyla yaklaşık \%2 ve \%5,5'e indirmiş,
PSNR'yi 55,5~dB'den 57,4~dB'ye yükseltmiştir. BT'nin girdi kanallarına eklenmesi (erken füzyon)
ise beklenenin aksine iyileşme sağlamamıştır: üç tohumun ortalamasında BT'siz model 21 hastanın
20'sinde daha yüksek PSNR vermiştir (Bonferroni düzeltmeli $p < 0{,}001$). Kübik
interpolasyonun $k = 4$'te lezyon SUVmax değerini ortalama \%34,5 düşük ölçmesi, problemin
seyrekleşme arttıkça ağırlaştığını göstermektedir. Bu bulgular hem hattın uygulanabilirliğini
hem de projenin ana sorusunu, yani BT bilgisinin ağa hangi yolla aktarılması gerektiğini
doğrudan desteklemektedir. Pilot sonuçlar Bölüm~\ref{sec:diger}'teki tabloda özetlenmiştir.

\Needspace{18\baselineskip}
\section{Proje Yönetimi}

\subsection{Çalışma Takvimi}

{\small
\begin{longtable}{@{}L{0.6cm}L{4.2cm}L{3.0cm}L{1.4cm}L{4.3cm}@{}}
\toprule
İP No & İş Paketinin Adı ve Hedefleri & Kim(ler) Tarafından Gerçekleştirileceği & Zaman Aralığı (Ay)
  & Başarı Ölçütü ve Projenin Başarısına Katkısı \\
\midrule
\endhead
1 & Veri hattı, kalite kontrolü ve bozunum: SUV dönüşümü, BT'nin PET ızgarasına örneklenmesi,
  hasta düzeyinde bölünme, $k = 2$, 3, 4 test bozunumlarının dondurulması, kübik taban çizgisi ve
  metrik altyapısı (Hedef 1)
  & Proje yürütücüsü, Proje ortağı 1 & 1--2
  & Kohortun tamamı için kalite kontrol raporu ve dışlama gerekçeleri; dondurulmuş bölünme ve
  bozunum dosyaları; kübik interpolasyonun üç faktördeki metrikleri. Katkı: \%15 \\
2 & Literatür yöntemlerinin uyarlanması: SAINT ve I$^3$Net'in ortalama tabanlı bozunumla yeniden
  eğitilmesi, RIFE ve FILM'in eğitimsiz ve ince ayarlı çalıştırılması (Hedef 2)
  & Proje ortağı 1, Proje ortağı 2 & 3--4
  & 4 yöntem $\times$ 3 faktörün ortak metrik koduyla değerlendirilmesi ve karşılaştırma
  tablosunun doldurulması. Katkı: \%15 \\
3 & Model merdiveni: BT'siz ve erken füzyonlu 2.5B U-Net ile BT temel modeli öznitelikli U-Net,
  her faktörde $\geq 3$ tohum (Hedef 3, H1)
  & Proje yürütücüsü, Proje ortağı 2 & 5--6
  & Öğrenmeli modelin kübik interpolasyona göre hasta düzeyinde anlamlı üstünlüğü (Bonferroni
  sonrası $p < 0{,}05$); BT'li ve BT'siz kollar arasındaki farkın test edilmesi. Katkı: \%20 \\
4 & Özgün çapraz dikkat modeli ve ablasyonlar: füzyon türü, girdi formülasyonu, kayıp
  bileşenleri, ölçüm tutarlılığı; opsiyonel ölçekten bağımsız sürüm (Hedef 4, H2, H3)
  & Proje yürütücüsü, Proje ortağı 1, Proje ortağı 2 & 7--10
  & $k = 4$'te BT'siz U-Net'e karşı PSNR'de anlamlı üstünlük; 1~mL altı lezyonlarda SUVmax
  düşük ölçümünün kübik interpolasyona göre en az \%40 azaltılması; ablasyon tablosu.
  Katkı: \%30 \\
5 & Genellenebilirlik ve klinik görev değerlendirmesi: farklı merkezden PSMA verisinde dış test,
  FDG ile tracer değişimi testi, segmentasyon tabanlı görev değerlendirmesi, hata vakası analizi
  (Hedef 5)
  & Proje ortağı 2, Proje yürütücüsü & 11--12
  & Dış test ve tracer değişimi sonuçlarının raporlanması; gerçek ve sentezlenmiş PET üzerindeki
  segmentasyon Dice farkı; en kötü beş vakanın analizi. Katkı: \%20 \\
\bottomrule
\end{longtable}
}

\Needspace{12\baselineskip}
\subsection{Risk Yönetimi}

{\small
\begin{longtable}{@{}L{1.0cm}L{5.0cm}L{8.0cm}@{}}
\toprule
İP No & En Önemli Riskler & Risk Yönetimi (B Planı) \\
\midrule
\endhead
1 & Yeni eklenen veri kümelerinde SUV dönüşümünün hatalı olması (bozunma düzeltmesi, enjekte doz
  veya kilo birimi)
  & Kalite kontrol kapısı, FDG'de karaciğer SUVmean kontrolü ve PSMA'da böbrek/mesane parlaklık
  sırası kontrolü uygulanır; sorunlu hastalar gerekçesiyle dışlanır ve kayda geçirilir.
  Hedeflerden sapma olmaz. \\
2 & Literatür yöntemlerinin açık kodlarının bu veriye veya ortalama tabanlı bozunuma
  uyarlanamaması
  & Yöntem, yayımlanmış mimarisine sadık kalınarak projenin kendi eğitim çerçevesinde yeniden
  gerçeklenir. Bir yöntem yine de çalıştırılamazsa nedeni raporlanır ve karşılaştırma kalan
  yöntemlerle yapılır; RIFE ve FILM'in eğitimsiz uygulanması her durumda mümkündür. \\
3, 4 & GPU hesaplama kotasının yetersiz kalması veya bulut oturumlarının kesilmesi
  & Tüm aşamalar kaldığı yerden sürdürülebilir tasarlanmıştır. Gerektiğinde karışık duyarlıklı
  eğitim, daha küçük yamalar ve temel modelin PET ızgarasındaki (daha ucuz) sürümü kullanılır;
  deneyler en bilgilendirici faktör olan $k = 4$'ten başlanarak sıralanır. Yöntem değişikliği
  gerekmez. \\
4 & Çapraz dikkat modelinin BT'siz modeli geçememesi
  & Bu sonuç araştırma sorusunun bilimsel bir cevabıdır ve hasta düzeyi istatistiklerle
  raporlanır. B planı olarak BT temel modelinin ince ayarı ve BT'nin yalnızca lezyon/mesane
  sınırları gibi seçilmiş bölgelerde kullanıldığı bir dikkat maskesi denenir; bu durumda yöntem
  değişikliği Bölüm~2.5'teki merdivene eklenen yeni bir basamak olarak tanımlanır ve temel
  hedefler değişmez. \\
4 & Yaklaşık 70 eğitim hastasıyla aşırı öğrenme
  & Rezidüel ve sıfırla başlatılmış tasarım, doğrulama kümesiyle model seçimi ve çoklu tohum
  uygulanır; gerekirse farklı merkezin PSMA verisinin bir bölümü eğitime eklenir ve dış test
  kalan hastalarla yapılır. \\
5 & Kurumsal veriye erişimin veya etik kurul onayının gecikmesi
  & Dış doğrulama, farklı bir merkezde çekilmiş açık PSMA-PET-CT-Lesions koleksiyonuyla yapılır;
  kurumsal veri yalnızca ek doğrulama olarak planlanmıştır. Hedeflerden sapma olmaz. \\
Tümü & Bir ekip üyesinin projeden ayrılması
  & Kod, yapılandırmalar ve deney günlüğü ortak bir sürüm kontrol deposunda tutulur; her iş
  paketinde en az iki kişi görevlidir. \\
\bottomrule
\end{longtable}
}

\subsection{Araştırma Olanakları}

{\small
\begin{tabularx}{\textwidth}{@{}YY@{}}
\toprule
Kuruluşta Bulunan Altyapı / Ekipman Türü & Projede Kullanım Amacı \\
\midrule
Hacettepe Üniversitesi Bilgisayar Mühendisliği Bölümü öğrenci laboratuvarları ve kampüs ağı
  & Geliştirme, kod incelemesi ve ekip toplantıları \\
{[Varsa: bölümün/fakültenin GPU sunucusu veya hesaplama kümesi]}
  & Çapraz dikkat modeli ve temel model deneyleri gibi yüksek bellek gerektiren eğitimler \\
Ekip üyelerine ait GPU'lu iş istasyonları (12--16~GB grafik belleği)
  & Kısa deneyler, hata ayıklama, çıkarım ve metrik hesapları \\
Bulut tabanlı GPU hesaplama ve bulut depolama hizmeti (proje bütçesinden)
  & Ana eğitim koşuları; ham ve işlenmiş verinin, kontrol noktalarının ve koşu arşivlerinin ekip
  arasında paylaşımlı saklanması \\
Açık kaynak yazılımlar (PyTorch, MONAI, SimpleITK, TotalSegmentator)
  & Model geliştirme, görüntü işleme ve organ maskeleri \\
\bottomrule
\end{tabularx}
}

\section{Araştırma Önerisinin Yaygın Etkisi}

{\small
\begin{tabularx}{\textwidth}{@{}L{4.5cm}Y@{}}
\toprule
Yaygın Etki Türü & Önerilen Araştırmadan Beklenen Çıktı, Sonuç ve Etkiler \\
\midrule
Bilimsel / Akademik Çıktılar (Ulusal/Uluslararası Makale, Bildiri, Kitap Bölümü vb.)
  & En az bir uluslararası tıbbi görüntüleme konferansına (ör. SPIE Medical Imaging, IEEE ISBI
  veya MICCAI çalıştayları) bildiri; nükleer tıp fiziği veya tıbbi fizik alanında SCI-E
  kapsamlı bir dergiye (ör. EJNMMI Physics, Medical Physics) makale taslağı; ulusal bir
  kongrede (ör. Ulusal Nükleer Tıp Kongresi veya SIU) Türkçe bildiri. PET kesit
  interpolasyonu için ortalama tabanlı bozunum ve lezyon düzeyi değerlendirme protokolü,
  sonraki çalışmaların kullanabileceği bir karşılaştırma standardı olarak sunulacaktır. \\
Ekonomik / Ticari / Sosyal Çıktılar (Ürün, Prototip, Yazılım, Veri Tabanı vb.)
  & Açık kaynak lisanslı yazılım: ön işleme ve kalite kontrol hattı, dondurulmuş bölünme ve
  bozunum protokolü, eğitilmiş model ağırlıkları ve değerlendirme betikleri. Uzun vadede
  yöntem, PET çekim süresini veya hasta dozunu artırmadan küçük lezyonların daha doğru
  ölçülmesine, bu sayede evreleme ve radyoligand tedavisine hasta seçiminde daha güvenilir
  nicel karar verilmesine katkı sağlayabilecektir. Kısa protokollerin hasta konforunu ve cihaz
  kapasitesini artırması dolaylı bir sosyal ve ekonomik etkidir. \\
Araştırmacı Yetiştirilmesi ve Yeni Proje(ler) Oluşturma
  & Üç lisans öğrencisinin tıbbi görüntü analizi, derin öğrenme ve tekrarlanabilir araştırma
  pratiklerinde deneyim kazanması; ekip üyelerinin lisansüstü çalışmalarına (yüksek lisans
  tezi) zemin oluşturması. Sonuçların, kurumsal PSMA verisiyle dış doğrulama ve klinik okuyucu
  çalışmasını içeren daha kapsamlı bir TÜBİTAK projesi (ör. 1001 veya 3501) başvurusunun ön
  verisi olarak kullanılması. \\
\bottomrule
\end{tabularx}
}

\Needspace{16\baselineskip}
\section{Belirtmek İstediğiniz Diğer Konular}\label{sec:diger}

Ön çalışma sonuçları ($k = 2$, DEEP-PSMA, 21 test hastası, 1.508 lezyon; ölçüm tutarlılığı
sonrası, ortalama $\pm$ standart sapma).

{\small
\begin{tabularx}{\textwidth}{@{}YL{2.6cm}L{3.2cm}L{3.6cm}@{}}
\toprule
Model & PSNR (dB) & Lezyon SUVmax sapması (\%) & 1~mL altı lezyonlarda SUVmax sapması (\%) \\
\midrule
Kübik spline & $55{,}46 \pm 7{,}44$ & $-13{,}2 \pm 11{,}1$ & $-15{,}7 \pm 11{,}0$ \\
2.5B U-Net, BT'siz (3 tohum) & 57,43 -- 57,53 & $-1{,}8$ ile $-2{,}3$ & $-5{,}2$ ile $-5{,}8$ \\
2.5B U-Net, BT erken füzyon (3 tohum) & 56,77 -- 57,18 & $-3{,}4$ ile $-3{,}5$ & $-7{,}1$ ile $-7{,}7$ \\
\bottomrule
\end{tabularx}
}

\paragraph{Etik.} Projede kullanılacak bütün veri setleri anonimleştirilmiş ve araştırma amaçlı
kullanım lisansıyla kamuya açık olarak yayımlanmıştır; hastalarla doğrudan temas veya yeni
görüntü çekimi yoktur. Bu nedenle üniversitenin etik kurulundan muafiyet belgesi alınması
planlanmaktadır. Kurumsal verinin kullanılması durumunda ilgili etik kurul onayı veri
kullanımından önce alınacaktır.

\paragraph{Üretken yapay zekâ kullanımı.} ``TÜBİTAK Destek Süreçlerinde Üretken Yapay Zekânın
Sorumlu ve Güvenilir Kullanımı Rehberi'' uyarınca beyan edilir: Proje metninin hazırlanmasında
üretken yapay zekâ araçlarından metin taslaklarının oluşturulması, ilgili literatürün taranması
ve kaynakların doğruluğunun (varlığı, künye bilgileri ve atıf yapılan içerikle uyumu) kontrol
edilmesi amacıyla yararlanılmıştır. Üretilen bütün içerik ekip üyeleri tarafından gözden
geçirilip düzenlenmiş ve akademik danışmanın kontrolüne sunulmuştur; kaynakların tamamı yayıncı
ve dizin kayıtları (Crossref, DataCite, arXiv, PubMed) üzerinden tek tek doğrulanmıştır.
Projenin bilimsel fikri, tasarımı ve içeriğinin sorumluluğu proje ekibine aittir.

\section{Ekler}

Bu başvuruda ek bulunmamaktadır.

\section*{Kaynaklar}

\begin{list}{}{\setlength{\leftmargin}{1.5em}\setlength{\itemindent}{-1.5em}%
  \setlength{\itemsep}{0.25\baselineskip}\setlength{\parsep}{0pt}}\small

\item Bland, J. M. ve Altman, D. G. 1986. ``Statistical methods for assessing agreement between two
methods of clinical measurement'', \emph{The Lancet}, 327(8476), 307--310.
\doi{10.1016/S0140-6736(86)90837-8}

\item Ben-Cohen, A., Klang, E., Raskin, S. P. vd. 2019. ``Cross-modality synthesis from CT to PET
using FCN and GAN networks for improved automated lesion detection'', \emph{Engineering
Applications of Artificial Intelligence}, 78, 186--194. \doi{10.1016/j.engappai.2018.11.013}

\item Boellaard, R., Delgado-Bolton, R., Oyen, W. J. G. vd. 2015. ``FDG PET/CT: EANM procedure
guidelines for tumour imaging: version 2.0'', \emph{European Journal of Nuclear Medicine and
Molecular Imaging}, 42(2), 328--354. \doi{10.1007/s00259-014-2961-x}

\item Buteau, J. P., Martin, A. J., Emmett, L. vd. 2022. ``PSMA and FDG-PET as predictive and
prognostic biomarkers in patients given [$^{177}$Lu]Lu-PSMA-617 versus cabazitaxel for metastatic
castration-resistant prostate cancer (TheraP): a biomarker analysis from a randomised, open-label,
phase 2 trial'', \emph{The Lancet Oncology}, 23(11), 1389--1397.
\doi{10.1016/S1470-2045(22)00605-2}

\item Cardoso, M. J., Li, W., Brown, R. vd. 2022. ``MONAI: An Open-Source Framework for Deep
Learning in Healthcare'', arXiv:2211.02701. \doi{10.48550/arXiv.2211.02701}

\item Chen, Y., Liu, S. ve Wang, X. 2021. ``Learning Continuous Image Representation with Local
Implicit Image Function'', \emph{2021 IEEE/CVF Conference on Computer Vision and Pattern
Recognition (CVPR)}, 8624--8634. \doi{10.1109/CVPR46437.2021.00852}

\item Clark, K., Vendt, B., Smith, K. vd. 2013. ``The Cancer Imaging Archive (TCIA): Maintaining
and Operating a Public Information Repository'', \emph{Journal of Digital Imaging}, 26(6),
1045--1057. \doi{10.1007/s10278-013-9622-7}

\item Dunn, O. J. 1961. ``Multiple Comparisons among Means'', \emph{Journal of the American
Statistical Association}, 56(293), 52--64. \doi{10.1080/01621459.1961.10482090}

\item Endo, H., Hirata, K., Magota, K. vd. 2025. ``Development and validation of 3D
super-resolution convolutional neural network for $^{18}$F-FDG-PET images'', \emph{EJNMMI Physics},
12(1), 77. \doi{10.1186/s40658-025-00791-y}

\item Fendler, W. P., Eiber, M., Beheshti, M. vd. 2023. ``PSMA PET/CT: joint EANM procedure
guideline/SNMMI procedure standard for prostate cancer imaging 2.0'', \emph{European Journal of
Nuclear Medicine and Molecular Imaging}, 50(5), 1466--1486. \doi{10.1007/s00259-022-06089-w}

\item Ferdinandus, J., Violet, J., Sandhu, S. vd. 2020. ``Prognostic biomarkers in men with
metastatic castration-resistant prostate cancer receiving [$^{177}$Lu]-PSMA-617'',
\emph{European Journal of Nuclear Medicine and Molecular Imaging}, 47(10), 2322--2327.
\doi{10.1007/s00259-020-04723-z}

\item Gafita, A., Wang, H., Robertson, A. vd. 2022. ``Tumor Sink Effect in $^{68}$Ga-PSMA-11 PET:
Myth or Reality?'', \emph{Journal of Nuclear Medicine}, 63(2), 226--232.
\doi{10.2967/jnumed.121.261906}

\item Gatidis, S., Hepp, T., Früh, M. vd. 2022. ``A whole-body FDG-PET/CT Dataset with manually
annotated Tumor Lesions'', \emph{Scientific Data}, 9(1), 601. \doi{10.1038/s41597-022-01718-3}

\item Hatamizadeh, A., Nath, V., Tang, Y. vd. 2022. ``Swin UNETR: Swin Transformers for Semantic
Segmentation of Brain Tumors in MRI Images'', \emph{Brainlesion: Glioma, Multiple Sclerosis,
Stroke and Traumatic Brain Injuries (BrainLes 2021)}, Lecture Notes in Computer Science, 12962,
272--284. \doi{10.1007/978-3-031-08999-2_22}

\item Hofman, M. S., Lawrentschuk, N., Francis, R. J. vd. 2020. ``Prostate-specific membrane
antigen PET-CT in patients with high-risk prostate cancer before curative-intent surgery or
radiotherapy (proPSMA): a prospective, randomised, multicentre study'', \emph{The Lancet},
395(10231), 1208--1216. \doi{10.1016/S0140-6736(20)30314-7}

\item Huang, Z., Zhang, T., Heng, W. vd. 2022. ``Real-Time Intermediate Flow Estimation for Video
Frame Interpolation'', \emph{Computer Vision -- ECCV 2022}, Lecture Notes in Computer Science,
13674, 624--642. \doi{10.1007/978-3-031-19781-9_36}

\item Isensee, F., Jaeger, P. F., Kohl, S. A. A. vd. 2021. ``nnU-Net: a self-configuring method
for deep learning-based biomedical image segmentation'', \emph{Nature Methods}, 18(2), 203--211.
\doi{10.1038/s41592-020-01008-z}

\item Jackson, P., Hofman, M., Buteau, J. P., McIntosh, L. ve Sun, Y. 2025. ``A Repository of
Annotated PSMA and FDG PET/CT Images for Algorithm Development in Staging of mCRPC for Treament
[sic] with 177Lu-PSMA Therapy'' [Veri seti], Zenodo. \doi{10.5281/zenodo.15281784}

\item Jeblick, K., Schachtner, B., Mittermeier, A. vd. 2026. ``A whole-body PSMA-PET/CT dataset
with manually annotated tumor lesions (PSMA-PET-CT-Lesions)'' (Sürüm 2) [Veri seti], The Cancer
Imaging Archive. \doi{10.7937/r7ep-3x37}

\item Koopman, D., van Dalen, J. A., Lagerweij, M. C. M. vd. 2015. ``Improving the Detection of
Small Lesions Using a State-of-the-Art Time-of-Flight PET/CT System and Small-Voxel
Reconstructions'', \emph{Journal of Nuclear Medicine Technology}, 43(1), 21--27.
\doi{10.2967/jnmt.114.147215}

\item Liu, J., Malekzadeh, M., Mirian, N. vd. 2021. ``Artificial Intelligence-Based Image
Enhancement in PET Imaging: Noise Reduction and Resolution Enhancement'', \emph{PET Clinics},
16(4), 553--576. \doi{10.1016/j.cpet.2021.06.005}

\item Loshchilov, I. ve Hutter, F. 2017. ``SGDR: Stochastic Gradient Descent with Warm Restarts'',
\emph{International Conference on Learning Representations (ICLR 2017)}. arXiv:1608.03983

\item Loshchilov, I. ve Hutter, F. 2019. ``Decoupled Weight Decay Regularization'',
\emph{International Conference on Learning Representations (ICLR 2019)}. arXiv:1711.05101

\item Pai, S., Hadzic, I., Bontempi, D. vd. 2025. ``Vision Foundation Models for Computed
Tomography'', arXiv:2501.09001. \doi{10.48550/arXiv.2501.09001}

\item Paszke, A., Gross, S., Massa, F. vd. 2019. ``PyTorch: An Imperative Style,
High-Performance Deep Learning Library'', \emph{Advances in Neural Information Processing Systems
32 (NeurIPS 2019)}.

\item Peng, C., Lin, W.-A., Liao, H. vd. 2020. ``SAINT: Spatially Aware Interpolation NeTwork for
Medical Slice Synthesis'', \emph{2020 IEEE/CVF Conference on Computer Vision and Pattern
Recognition (CVPR)}, 7747--7756. \doi{10.1109/CVPR42600.2020.00777}

\item Reda, F., Kontkanen, J., Tabellion, E. vd. 2022. ``FILM: Frame Interpolation for Large
Motion'', \emph{Computer Vision -- ECCV 2022}, Lecture Notes in Computer Science, 13667,
250--266. \doi{10.1007/978-3-031-20071-7_15}

\item Ronneberger, O., Fischer, P. ve Brox, T. 2015. ``U-Net: Convolutional Networks for
Biomedical Image Segmentation'', \emph{Medical Image Computing and Computer-Assisted Intervention
-- MICCAI 2015}, Lecture Notes in Computer Science, 9351, 234--241.
\doi{10.1007/978-3-319-24574-4_28}

\item Sanayi ve Teknoloji Bakanlığı. 2025. ``2030 Sanayi ve Teknoloji Stratejisi''.
\url{https://www.sanayi.gov.tr/assets/pdf/plan-program/2030SanayiveTeknolojiStratejisi.pdf}
Son erişim tarihi: 15 Kasım 2026.

\item Sartor, O., de Bono, J., Chi, K. N. vd. 2021. ``Lutetium-177--PSMA-617 for Metastatic
Castration-Resistant Prostate Cancer'', \emph{New England Journal of Medicine}, 385(12),
1091--1103. \doi{10.1056/NEJMoa2107322}

\item Schlemper, J., Caballero, J., Hajnal, J. V. vd. 2018. ``A Deep Cascade of Convolutional
Neural Networks for Dynamic MR Image Reconstruction'', \emph{IEEE Transactions on Medical
Imaging}, 37(2), 491--503. \doi{10.1109/TMI.2017.2760978}

\item Song, H., Mao, X., Yu, J. vd. 2024. ``I$^3$Net: Inter-Intra-Slice Interpolation Network for
Medical Slice Synthesis'', \emph{IEEE Transactions on Medical Imaging}, 43(9), 3306--3318.
\doi{10.1109/TMI.2024.3394033}

\item Song, T.-A., Chowdhury, S. R., Yang, F. ve Dutta, J. 2020. ``Super-Resolution PET Imaging
Using Convolutional Neural Networks'', \emph{IEEE Transactions on Computational Imaging}, 6,
518--528. \doi{10.1109/TCI.2020.2964229}

\item Soret, M., Bacharach, S. L. ve Buvat, I. 2007. ``Partial-Volume Effect in PET Tumor
Imaging'', \emph{Journal of Nuclear Medicine}, 48(6), 932--945. \doi{10.2967/jnumed.106.035774}

\item Strateji ve Bütçe Başkanlığı. 2023. ``On İkinci Kalkınma Planı (2024--2028)''.
\url{https://www.sbb.gov.tr/wp-content/uploads/2023/12/On-Ikinci-Kalkinma-Plani_2024-2028_11122023.pdf}
Son erişim tarihi: 15 Kasım 2026.

\item TÜBİTAK. 2025. ``Öncelikli Alanlar''.
\url{https://tubitak.gov.tr/sites/default/files/2025-07/Oncelikli_Alanlar_2025-2.pdf}
Son erişim tarihi: 15 Kasım 2026.

\item Vaswani, A., Shazeer, N., Parmar, N. vd. 2017. ``Attention Is All You Need'', \emph{Advances
in Neural Information Processing Systems 30 (NIPS 2017)}. arXiv:1706.03762

\item Wang, Z., Bovik, A. C., Sheikh, H. R. ve Simoncelli, E. P. 2004. ``Image Quality Assessment:
From Error Visibility to Structural Similarity'', \emph{IEEE Transactions on Image Processing},
13(4), 600--612. \doi{10.1109/TIP.2003.819861}

\item Wang, Z., Simoncelli, E. P. ve Bovik, A. C. 2003. ``Multiscale Structural Similarity for
Image Quality Assessment'', \emph{The Thirty-Seventh Asilomar Conference on Signals, Systems \&
Computers}, 1398--1402. \doi{10.1109/ACSSC.2003.1292216}

\item Wasserthal, J., Breit, H.-C., Meyer, M. T. vd. 2023. ``TotalSegmentator: Robust
Segmentation of 104 Anatomic Structures in CT Images'', \emph{Radiology: Artificial
Intelligence}, 5(5), e230024. \doi{10.1148/ryai.230024}

\item Wilcoxon, F. 1945. ``Individual Comparisons by Ranking Methods'', \emph{Biometrics
Bulletin}, 1(6), 80--83. \doi{10.2307/3001968}

\item Wu, Q., Li, Y., Sun, Y. vd. 2023. ``An Arbitrary Scale Super-Resolution Approach for 3D MR
Images via Implicit Neural Representation'', \emph{IEEE Journal of Biomedical and Health
Informatics}, 27(2), 1004--1015. \doi{10.1109/JBHI.2022.3223106}

\item Zhao, C., Dewey, B. E., Pham, D. L. vd. 2021. ``SMORE: A Self-Supervised Anti-Aliasing and
Super-Resolution Algorithm for MRI Using Deep Learning'', \emph{IEEE Transactions on Medical
Imaging}, 40(3), 805--817. \doi{10.1109/TMI.2020.3037187}

\item Zhao, H., Gallo, O., Frosio, I. ve Kautz, J. 2017. ``Loss Functions for Image Restoration
With Neural Networks'', \emph{IEEE Transactions on Computational Imaging}, 3(1), 47--57.
\doi{10.1109/TCI.2016.2644865}

\end{list}

\section*{Bütçe Talebi (TYBS bütçe adımı için öneri)}

{\small
\begin{tabularx}{\textwidth}{@{}L{2.6cm}L{3.6cm}Yr@{}}
\toprule
Gider Türü & Kalem & Gerekçe & Tutar (TL) \\
\midrule
Hizmet alımı & Bulut tabanlı GPU hesaplama hizmeti (12 ay)
  & İP3--İP5'teki eğitim koşuları; çapraz dikkat ve temel model deneyleri yüksek bellekli GPU
  gerektirir & [7.000] \\
Hizmet alımı & Bulut depolama hizmeti (2~TB, 12 ay)
  & Yaklaşık 40~GB ham veri, işlenmiş ve bozunmuş sürümler, kontrol noktaları ve koşu
  arşivlerinin ekip arasında paylaşımı & [2.500] \\
Sarf malzemesi & Harici SSD (2~TB)
  & Verinin ve sonuçların yerel yedeği, bulut kesintilerinde yerel iş istasyonunda çalışma
  & [2.500] \\
\midrule
 & & Toplam & [12.000] \\
\bottomrule
\end{tabularx}
}

\end{document}
